\section{Discussão}
\label{sec:discussion}

\subsection{O que a replicação mostra sobre o \modus}
A replicação separa três afirmações que o artigo original apresenta juntas. A primeira é que o
método atinge os F1 reportados. Isso se reproduz em duas das três bases, mas no ASEADOS só com o
rótulo no fluxo, e no Edge-IIoTset com VP e FN idênticos em um conjunto que conta o tráfego normal
duas vezes. A segunda é que a detecção é volumétrica. Não é: sem normalização, o HST responde à
raridade do protocolo, e o recall de 73,16\% no Edge-IIoTset é exatamente a fração de pacotes DDoS
que não são TCP. A terceira é que a \emph{flag} de mudança de conceito informa o detector. Ela muda
um pacote em 19,6 milhões.

Não afirmamos que os autores usaram o rótulo. O código não é público, e o texto não diz como o
fluxo foi montado. O que os dados sustentam é mais restrito: com o HST do MOA e o uso padrão do
CapyMOA, o rótulo chega ao detector, e o número do ASEADOS se reproduz com ele e não se
reproduz sem ele. A coincidência justifica pedir aos autores a publicação do código.

\subsection{Um vazamento que a API não mostra}
O vazamento descrito na Seção~\ref{sec:leak} não exige erro do usuário. O \texttt{NumpyStream} do
CapyMOA recebe \texttt{X} e \texttt{y}, como qualquer fluxo de avaliação, e o \texttt{HSTrees}
aparece na biblioteca como detector de anomalia de uma classe. Nada na interface indica que o
atributo de classe entra nas dimensões de corte. O efeito também não é uniforme, e isso o torna
difícil de notar: no CTU-13 as predições com e sem rótulo são idênticas, porque os atributos já
separam as classes. Um teste em uma só base poderia concluir que não há problema.

Duas verificações simples teriam revelado o defeito, e as recomendamos para qualquer avaliação de
detector não supervisionado: (i) repetir a rodada com rótulos constantes ou embaralhados, pois o
resultado não pode mudar; (ii) pontuar o mesmo vetor de atributos com rótulos diferentes. Ambas se
enquadram no que Arp \emph{et al.}~\cite{arp2022dos} chamam de checagem de \emph{label leakage}, e
custam uma rodada a mais.

\subsection{Envenenamento sem adversário}
A literatura de envenenamento costuma supor um atacante que conhece o detector e injeta tráfego
aos poucos~\cite{rubinstein2009antidote}. Nossos resultados mostram um caso mais simples e mais
comum: o próprio ataque volumétrico, sem nenhuma intenção de evasão, envenena o modelo. Bastam
duas condições, ambas presentes no \modus: o detector treina com o que classifica como benigno, e
o ritmo de atualização do modelo é medido em instâncias. Um \emph{flood} produz instâncias na
taxa que quiser; logo, controla o relógio do modelo.

Daí a utilidade de reportar a fração do ataque usada no treino. O F1 diz se o detector acertou; a
fração treinada diz se ele continuará acertando. No CTU-13, por exemplo, o modelo congelado
sozinho tem o melhor F1 a priori (69,59\%), mas 43,58\% do \emph{flood} treina o modelo adaptativo.
Já a combinação de janela em tempo e teto tem F1 nulo e um modelo íntegro, com AUC de 96,47\%. Só
o F1 esconderia essa diferença.

\subsection{O papel de cada defesa}
As quatro defesas agem em pontos diferentes, e as bases as exercitam de modo diferente. A janela em
tempo e o teto limitam a \emph{velocidade} com que o ataque entra no modelo. São decisivos no
CTU-13, onde o ataque chega a 20 mil pacotes por segundo, e quase irrelevantes no ASEADOS, onde o
DDoS tem cerca de 46 fluxos por segundo e o teto de 24,9 por destino quase não atua: com e sem
ele, o envenenamento é de 51,24\% e 52,93\%. O modelo congelado age sobre a \emph{decisão}: ele mantém uma referência que o ataque não
alcança. É decisivo no ASEADOS e necessário em todas as bases, pois sem ele o limiar adaptativo
não dispara mesmo quando o score separa bem as classes. O aquecimento adaptativo decide \emph{com
que dados} o congelado é construído, e seu efeito no ASEADOS foi maior do que esperávamos.

O modelo congelado tem um custo conhecido, e os dados o confirmam. Ele pressupõe aquecimento limpo:
no ASEADOS com cinco classes, em que há ataques desde o início do fluxo, sua AUC é 50,05\%. E ele não
acompanha a deriva do tráfego legítimo: os falsos positivos sobem para 3,14\% no Edge-IIoTset e
17,88\% no ASEADOS com cinco classes. Em operação real, um modelo congelado precisaria ser refeito
periodicamente, em janelas validadas como limpas, o que reintroduz alguma supervisão.

\subsection{O limiar é o gargalo}
O resultado mais consistente entre as bases é a distância entre separabilidade e decisão. No
ASEADOS o score congelado tem AUC de 94,44\% e o recall é 47,07\%. No CTU-13 a AUC é 93,95\% e o
recall é 0,00\%. Nos dois casos o modelo ordena bem as instâncias, e a regra de limiar, um quantil
alto dos scores do aquecimento, corta no lugar errado. Um quantil estima a cauda do tráfego
\emph{benigno}; ele nada diz sobre onde o ataque vai cair. Com 1.331 a 7.064 instâncias de
calibração, o quantil 0,999 é determinado por uma a sete observações, o que também explica a
variação de $\tau_f$ entre sementes.

Vemos três caminhos, todos sem rótulos: (i) decidir por janela, e não por instância, exigindo que
uma fração dos scores de um destino passe de um quantil mais baixo, o que explora a persistência do
\emph{flood}; (ii) combinar os scores dos dois modelos em vez de fazer um \emph{ou} lógico entre
dois limiares; (iii) calibrar o limiar com a teoria de valores extremos sobre a cauda do
aquecimento, em vez de um quantil empírico. Não testamos nenhum deles.

\subsection{Ameaças à validade}
\textbf{Ordem do ASEADOS.} A ordem ``temporal'' que usamos é a que reproduz o artigo, e ela vem de
uma leitura errada das datas. Na ordem cronológica correta, o DDoS ocupa os primeiros 50 mil
fluxos, e nenhum método que aprende o normal no início do fluxo pode detectá-lo. Os resultados do
ASEADOS valem para a ordem do artigo, não para a base como foi coletada.

\textbf{Rótulos por pacote.} No CTU-13 o rótulo vem do IP infectado, e no Edge-IIoTset, do arquivo
de captura. Pacotes legítimos de uma máquina infectada, ou tráfego de fundo dentro de um
\emph{pcap} de ataque, contam como ataque. Isso limita o recall alcançável e afeta todos os métodos
igualmente.

\textbf{Uma rodada por base.} Só o ASEADOS foi rodado com três sementes. No CTU-13 e no
Edge-IIoTset há uma rodada por configuração. Dada a variação de dez pontos entre sementes no
ASEADOS, diferenças pequenas entre variantes não devem ser interpretadas.

\textbf{Parâmetros a priori.} A configuração \texttt{fixed\_q} e os parâmetros novos foram fixados
antes das rodadas da Etapa~2, mas por nós, que já conhecíamos as bases da Etapa~1. As rodadas de
sensibilidade estão identificadas como tal, e o resultado negativo do CTU-13 foi mantido como
resultado principal para não esconder essa dependência.

\textbf{Bases de laboratório.} As três bases têm ataques longos, intensos e de vetor único, depois
de um período benigno. Ataques curtos, de baixa taxa ou presentes desde o início não estão
representados.
